\documentclass[12pt, a4paper]{article}

% -----------------------------------------------------------
% Packages
% -----------------------------------------------------------

\usepackage[utf8]{inputenc} % UTF-8 encoding
\usepackage[margin=1in]{geometry} % Page geometry
\usepackage{amsmath, amssymb} % Math symbols and equations
\usepackage{graphicx} % Include images
\usepackage{hyperref} % Clickable links in PDF
\usepackage{setspace} % Line spacing
\usepackage{xcolor} % Text color
\usepackage{titlesec} % Custom section titles
\usepackage{enumitem} % Custom lists

% -----------------------------------------------------------
% Document Configuration
% -----------------------------------------------------------

\title{Metodi Formali}
\author{Claudio Tessa}
\date{\the\year}
    
% Custom section styling
\titleformat{\section}
  {\normalfont\large\bfseries}{\thesection}{1em}{}

% -----------------------------------------------------------
% Begin Document
% -----------------------------------------------------------

\begin{document}
\maketitle

\section*{Informazioni generali}
L'esame consiste in uno scritto e in un orale. Il giorno dell'orale verrà concordato al momento della prova scritta. La prova orale consiste nella discussione dello scritto e inoltre verranno fatte domande sugli argomenti trattati nel corso.

\tableofcontents
\newpage

% -----------------------------------------------------------
% Introduzione
% -----------------------------------------------------------

\section{Introduzione}

Verranno trattati metodi e tecniche formali per specificare, disegnare e analizzare sistemi complessi. In particolare, \textbf{sistemi concorrenti e distribuiti}, costituiti da componenti che operano in modo indipendente e che interagiscono tra loro. Ci concentreremo sulle fasi di \textbf{specifica del problema} e \textbf{modellazione di una soluzione}. L'uso di metodi formali in queste fasi permette di verificare e validare il modello prima di procedere all'implementazione.
\\
\\
Per fare ciò verrano introdotti:
\begin{itemize}
    \item Un \textbf{linguaggio logico} per specificare le proprietà di comportamento di tali sistemi
    \item Le \textbf{reti di Petri} come strumento per modellare i sistemi concorrenti e analizzarne le proprietà    
    \item Nozioni di base sui sistemi dinamici a tempi discreti (cenni ad automi cellulari)
\end{itemize}

\subsection{Teoria Generale delle Reti di Petri}

L'obiettivo di Petri era sviluppare una teoria matematica fondata sui principi della fisica moderna (teoria della relatività e fisica quantistica), che sia una \textbf{teoria dei sistemi} in grado di descrivere il \textbf{flusso di informazione} e permetta di analizzare \textbf{sistemi con organizzazione complessa}. Una teoria generale dei sistemi che processano informazione basata su: comunicazione, sincronizzazione, flusso di informazione, etc.

% -----------------------------------------------------------
% End Document
% -----------------------------------------------------------

\end{document}
